\documentclass[12pt]{report}

\usepackage[a4paper,margin=1in]{geometry}
\usepackage[T1]{fontenc}
\usepackage[utf8]{inputenc}
\usepackage{lmodern}
\usepackage{graphicx}
\usepackage{booktabs}
\usepackage{tabularx}
\usepackage{array}
\usepackage{ragged2e}
\usepackage{enumitem}
\usepackage{listings}
\usepackage{xcolor}
\usepackage{float}
\usepackage[hidelinks]{hyperref}
\usepackage{caption}

\captionsetup{font=small}
\setlist[itemize]{leftmargin=*,itemsep=0.35em}
\setlist[enumerate]{leftmargin=*,itemsep=0.35em}
\setlength{\parskip}{0.6em}
\setlength{\parindent}{0pt}
\renewcommand{\arraystretch}{1.25}

\definecolor{codebg}{RGB}{248,248,248}
\definecolor{placeholdergray}{RGB}{110,110,110}

\lstdefinestyle{reportcode}{
  basicstyle=\ttfamily\small,
  backgroundcolor=\color{codebg},
  frame=single,
  breaklines=true,
  columns=fullflexible,
  keepspaces=true
}

\newcolumntype{Y}{>{\RaggedRight\arraybackslash}X}
\newcommand{\placeholder}[1]{\textcolor{placeholdergray}{\textit{#1}}}
\newcommand{\evidencebox}[2]{
  \begin{figure}[H]
    \centering
    \includegraphics[width=0.88\textwidth]{#1}
    \caption{#2}
  \end{figure}
}

\newcommand{\evidenceplaceholder}[2]{
  \begin{figure}[H]
    \centering
    \fbox{
      \parbox{0.88\textwidth}{
        \centering
        \vspace{1.2in}
        \placeholder{#1}
        \vspace{1.2in}
      }
    }
    \caption{#2}
  \end{figure}
}

\newcommand{\lessonchapter}[3]{
  \chapter{Lesson #1: #2}

  \section{Part 1: Goal and Vulnerability Summary}
  \placeholder{State the vulnerability type, main affected component, security impact, and the expected safe behavior for #2.}

  \section{Part 2: Why This Works / Root Cause}
  \placeholder{Explain the security mistake that makes this #3 issue possible. Keep this conceptual and focused on root cause rather than symptoms.}

  \section{Part 3: Environment and Setup}
  \begin{itemize}
    \item DVSA Website URL / API endpoint: https://shfz8h7a28.execute-api.us-east-1.amazonaws.com/Stage
    \item AWS services involved: \placeholder{For example API Gateway, Lambda, Cognito, S3, DynamoDB, SQS, IAM.}
    \item Relevant function, role, or log group: \placeholder{Name the exact resources used in this lesson.}
    \item User or test context: \placeholder{Attacker, victim, admin, or anonymous context as applicable.}
    \item Tools used: \placeholder{Browser DevTools, curl, jq, AWS Console, AWS CLI, Policy Simulator, CloudWatch, etc.}
  \end{itemize}

  \section{Part 4: Reproduction Steps}
  \begin{enumerate}
    \item \placeholder{Describe the baseline or normal behavior first.}
    \item \placeholder{List the exact steps that trigger the vulnerability.}
    \item \placeholder{Include the request, timing, role, or payload details needed for reproducibility.}
    \item \placeholder{Record the observable result that proves the issue.}
  \end{enumerate}

  \section{Part 5: Evidence and Proof}
  \begin{itemize}
    \item \placeholder{Attach the strongest pre-fix evidence for this lesson.}
    \item \placeholder{Use only evidence that directly proves the claim.}
    \item \placeholder{Redact real tokens, secrets, and temporary credentials before submission.}
  \end{itemize}
  \evidenceplaceholder{Replace with pre-fix evidence for Lesson #1 (#2).}{Pre-fix evidence for Lesson #1: #2}

  \section{Part 6: Fix Strategy / Probable Mitigation}
  \begin{itemize}
    \item Fix location: \placeholder{Identify the backend, IAM, cloud configuration, or workflow layer that must change.}
    \item Security property restored: \placeholder{For example authentication integrity, authorization, least privilege, or safe input handling.}
    \item Why this addresses root cause: \placeholder{Explain why the fix is real and not cosmetic.}
  \end{itemize}

  \section{Part 7: Code / Config Changes}
  \placeholder{Paste the real code diff, policy change, CloudFormation change, or configuration update below.}

  \begin{figure}[H]
    \centering
    \fbox{
      \parbox{0.88\textwidth}{
        \vspace{0.4em}
        \placeholder{Replace this box with the actual diff, code snippet, policy JSON, or configuration change.}
        \vspace{2.0in}
      }
    }
    \caption{Replace with the concrete code or configuration change for Lesson #1.}
  \end{figure}

  \section{Part 8: Verification After Fix}
  \begin{enumerate}
    \item \placeholder{Repeat the same test that proved the vulnerability.}
    \item \placeholder{Show that the exploit no longer succeeds.}
    \item \placeholder{Also show that normal intended behavior still works.}
  \end{enumerate}
  \evidenceplaceholder{Replace with post-fix verification evidence for Lesson #1 (#2).}{Post-fix verification evidence for Lesson #1: #2}

  \section[Part 9: Structured Operation and Security Analysis]{Part 9: Structured Operation and\\Security Analysis}
  \subsection*{Table A}
  \begin{table}[H]
    \small
    \centering
    \begin{tabularx}{\textwidth}{|Y|Y|Y|Y|}
      \hline
      \textbf{Intended rule(s)} & \textbf{Artifacts used to infer the rule} & \textbf{Normal evidence} & \textbf{Exploit evidence} \\
      \hline
      \placeholder{State the intended security or workflow rule.} &
      \placeholder{Identify the source used to infer the rule: UI flow, code path, IAM policy, logs, helper guide, etc.} &
      \placeholder{Show normal behavior when the rule is respected.} &
      \placeholder{Show behavior proving the rule was violated.} \\
      \hline
    \end{tabularx}
  \end{table}

  \subsection*{Table B}
  \begin{table}[H]
    \footnotesize
    \centering
    \begin{tabularx}{\textwidth}{|Y|Y|Y|Y|Y|}
      \hline
      \textbf{Why the exploit is a deviation} & \textbf{Deviation class} & \textbf{Fix applied} & \textbf{Post-fix verification} & \textbf{Optional latency / timing note} \\
      \hline
      \placeholder{Explain why the observed result violates the intended rule.} &
      \placeholder{Examples: authentication failure, access control failure, unsafe input handling, logic flaw, DoS, excessive privilege.} &
      \placeholder{Summarize the real corrective action.} &
      \placeholder{Describe the evidence that the deviation is closed.} &
      \placeholder{Use this only when timing or latency matters.} \\
      \hline
    \end{tabularx}
  \end{table}

  \section{Part 10: Takeaway / Lessons Learned}
  \placeholder{End with a short security lesson that ties the #3 weakness back to serverless design, AWS service interaction, and secure engineering practice.}

  \clearpage
}

% Fill in these fields before compiling the final report.
\newcommand{\CourseCode}{ICS-344}
\newcommand{\CourseTitle}{Information Security}
\newcommand{\CourseSection}{Sections: 1 \& 2}
\newcommand{\CourseTerm}{252}
\newcommand{\ProjectTitle}{DVSA Vulnerability Discovery and Remediation}
\newcommand{\StudentNames}{Husain Al Muallim$^1$, Ali Al Qatari$^2$, Abdulwahed Zuhair Adi$^3$}
\newcommand{\StudentIDs}{202163730, 202279780, 202256380}
\newcommand{\DVSAWebsiteURL}{http://dvsa-website-145292399251-us-east-1.s3-website.us-east-1.amazonaws.com}
\newcommand{\AWSRegion}{us-east-1}
\newcommand{\SubmissionDate}{\placeholder{Submission date}}

\begin{document}

\hypersetup{pageanchor=false}
\begin{titlepage}
  \centering
  {\Large King Fahd University of Petroleum and Minerals\par}
  \vspace{1.2cm}
  {\Large \CourseCode\ - \CourseTitle\par}
  {\large \CourseSection\par}
  \vspace{1.4cm}
  {\huge \bfseries \ProjectTitle\par}
  \vspace{1.4cm}

  \begin{tabularx}{\textwidth}{@{}rY@{}}
    Term: & \CourseTerm \\
    Student Name(s): & \StudentNames \\
    Student ID(s): & \StudentIDs \\
    AWS Region: & \AWSRegion \\
    Date: & \SubmissionDate \\
  \end{tabularx}

  \vspace{0.6cm}
  {\small\textbf{DVSA Website URL:} \url{\DVSAWebsiteURL}\par}

  \vfill
  {\large Project Report\par}
\end{titlepage}

\clearpage
\hypersetup{pageanchor=true}
\pagenumbering{roman}
\tableofcontents
\clearpage

\chapter*{Project Overview}
\addcontentsline{toc}{chapter}{Project Overview}

\section*{Scope}
\placeholder{Summarize the project scope: deployment in a non-production AWS account, the 10 official lessons, remediation, and verification.}

\section*{Deployment Summary}
\begin{itemize}
  \item DVSA Website URL: \DVSAWebsiteURL
  \item AWS region: \AWSRegion
  \item CloudFormation stack: \placeholder{Stack name}
  \item Main services used: \placeholder{S3, API Gateway, Lambda, DynamoDB, SQS, Cognito, IAM, CloudWatch}
\end{itemize}

\section*{Architecture and Request Pipeline}
\placeholder{Explain the intended high-level request flow from browser to API Gateway to Lambda to backend services and back.}
\evidenceplaceholder{Optional architecture figure or request-flow diagram.}{Architecture figure or request-flow diagram}

\section*{Methodology and Evidence Sources}
\begin{itemize}
  \item \placeholder{How the environment was validated before testing.}
  \item \placeholder{What tools were used to reproduce and verify vulnerabilities.}
  \item \placeholder{How evidence was captured and redacted.}
  \item \placeholder{How fixes were verified after remediation.}
\end{itemize}

\clearpage
\pagenumbering{arabic}

\chapter{Lesson 01: Event Injection}

\section{Part 1: Goal and Vulnerability Summary}
The objective of this lesson is to successfully exploit an Event Injection flaw to achieve Remote Code Execution (RCE) on the backend cloud server, and subsequently remediate the vulnerability. The vulnerability lies within the \texttt{DVSA-ORDER-MANAGER} AWS Lambda function, which acts as the backend for the application's order system. When a user submits an HTTP request, the API Gateway passes the user's JSON payload directly to this Lambda function. Instead of safely parsing this text, the developer used an outdated and inherently insecure Node.js library called \texttt{node-serialize} to deserialize both the request body and the HTTP headers. Because the backend blindly trusts the incoming user input and processes it with this vulnerable library, an attacker can craft a malicious JSON payload containing executable code, bypassing normal application logic.

\section{Part 2: Why This Works / Root Cause}
The root cause is insecure deserialization. The \texttt{DVSA-ORDER-MANAGER} Lambda function processes incoming HTTP requests using \texttt{serialize.unserialize(event.body)}. Under the hood, the \texttt{node-serialize} package uses JavaScript's \texttt{eval()} function to execute code if it detects a specifically formatted string (starting with \texttt{\_\$\$ND\_FUNC\$\$\_}). Because the backend does not validate or sanitize the incoming JSON payload before passing it to this function, the attacker's injected JavaScript payload is executed with the privileges of the Lambda function.

\section{Part 3: Environment and Setup}
\begin{itemize}
  \item API endpoint:\par
  \url{https://shfz8h7a28.execute-api.us-east-1.amazonaws.com/Stage}
  \item AWS services involved: API Gateway, AWS Lambda, CloudWatch Logs.
  \item Relevant function, role, or log group: Lambda function \nolinkurl{DVSA-ORDER-MANAGER} and its CloudWatch Log Group.
  \item User or test context: Unauthenticated external attacker.
  \item Tools used: Mac Terminal, \texttt{curl} for sending HTTP POST requests, and AWS Management Console (CloudWatch).
\end{itemize}

\section{Part 4: Reproduction Steps}
\begin{enumerate}
  \item The baseline behavior of the \texttt{/order} endpoint is to accept a valid JSON payload containing order details and process it.
  \item To trigger the vulnerability, a crafted JSON payload is sent to the \texttt{/order} API endpoint using \texttt{curl}.
  \item The payload contains the malicious \texttt{\_\$\$ND\_FUNC\$\$\_} tag, followed by a self-executing JavaScript function that writes a file to the Lambda's \texttt{/tmp} directory, reads it, and prints the contents to the console error log.
  \item The following command was executed:
\begin{lstlisting}[style=reportcode]
curl -X POST "https://shfz8h7a28.execute-api.us-east-1.amazonaws.com/Stage/order" \
 -H "Content-Type: application/json" \
 -d '{"action": "_$$ND_FUNC$$_function(){ var fs = require(\"fs\"); fs.writeFileSync(\"/tmp/pwned.txt\", \"You are reading the contents of my hacked file!\"); var fileData = fs.readFileSync(\"/tmp/pwned.txt\", \"utf-8\"); console.error(\"FILE READ SUCCESS: \" + fileData); }()", "cart-id":""}'
\end{lstlisting}
  \item The terminal returns an \texttt{\{"message": "Internal server error"\}}, indicating the payload executed and crashed the backend function.
\end{enumerate}

\section{Part 5: Evidence and Proof}
\begin{itemize}
  \item CloudWatch logs were inspected to verify that the payload executed successfully.
  \item The logs clearly show the injected message "FILE READ SUCCESS: You are reading the contents of my hacked file!", proving that the remote code execution was successful.
\end{itemize}
\evidencebox{../evidence/lesson-01/terminal-output.png}{Screenshot of Terminal Output showing the Internal Server Error}
\evidencebox{../evidence/lesson-01/cloudwatch-log.png}{Screenshot of CloudWatch Log showing the "FILE READ SUCCESS" message}


\section{Part 6: Fix Strategy / Probable Mitigation}
\begin{itemize}
  \item Fix location: Backend Lambda \nolinkurl{DVSA-ORDER-MANAGER}, file \nolinkurl{order-manager.js}.
  \item Security property restored: Safe input handling and prevention of Remote Code Execution.
  \item Why this addresses root cause: The mitigation strategy requires removing the insecure deserialization mechanism entirely. The \texttt{node-serialize} library must be removed from the code. All incoming JSON data (both the event body and headers) must be parsed using the native, safe \texttt{JSON.parse()} method, which only interprets data as text/objects and never executes it as code.
\end{itemize}

\section{Part 7: Code / Config Changes}
The \nolinkurl{order-manager.js} file in the \nolinkurl{DVSA-ORDER-MANAGER} Lambda function was updated to replace the insecure library.

The following lines were modified:
\begin{itemize}
  \item \textbf{Line 1:} \texttt{// const serialize = require('node-serialize');} (Commented out)
  \item \textbf{Line 10:} \texttt{var req = JSON.parse(event.body);} (Replaced unserialize)
  \item \textbf{Line 11:} \texttt{var headers = JSON.parse(event.headers);} (Replaced unserialize)
\end{itemize}

\evidenceplaceholder{Insert screenshot of your corrected Lambda code showing JSON.parse here.}{Corrected Lambda code for Lesson 01}

% \section{Part 8: Verification After Fix}
\begin{enumerate}
  \item After deploying the updated Lambda function, the exact same malicious \texttt{curl} command from Part 4 was executed in the terminal.
  \item The terminal again returned an \texttt{\{"message": "Internal server error"\}} (due to a missing authentication header failing later in the safe code).
  \item The CloudWatch logs for the \texttt{DVSA-ORDER-MANAGER} were reviewed again. The attack failed to execute, and the \texttt{FILE READ SUCCESS} message was no longer present in the logs, confirming the vulnerability is patched.
\end{enumerate}
\evidenceplaceholder{Insert screenshot of the clean CloudWatch logs here, showing no successful hack.}{Post-fix CloudWatch logs for Lesson 01}

% \section{Part 9: Structured Operation and Security Analysis}
\subsection*{Table A: Operational Analysis}
\begin{table}[H]
  \small
  \centering
  \begin{tabularx}{\textwidth}{|Y|Y|Y|Y|}
    \hline
    \textbf{Intended rule(s)} & \textbf{Artifacts used to infer the rule} & \textbf{Normal evidence} & \textbf{Exploit evidence} \\
    \hline
    The backend should safely parse JSON data representing an order without executing arbitrary code contained within the payload. &
    Source code analysis of \nolinkurl{order-manager.js} and normal application workflow. &
    The API accepts standard JSON payloads and processes order updates. &
    CloudWatch logs display output from custom JavaScript injected by the attacker. \\
    \hline
  \end{tabularx}
\end{table}

\subsection*{Table B: Security Analysis}
\begin{table}[H]
  \footnotesize
  \centering
  \begin{tabularx}{\textwidth}{|Y|Y|Y|Y|Y|}
    \hline
    \textbf{Why the exploit is a deviation} & \textbf{Deviation class} & \textbf{Fix applied} & \textbf{Post-fix verification} & \textbf{Optional latency / timing note} \\
    \hline
    The application evaluated user input as code rather than pure data. &
    Unsafe input handling / Insecure Deserialization. &
    Removed the \nolinkurl{node-serialize} module and replaced \nolinkurl{serialize.unserialize()} calls with \nolinkurl{JSON.parse()}. &
    Re-running the exploit payload no longer prints the injected message to the CloudWatch logs. &
    N/A \\
    \hline
  \end{tabularx}
\end{table}

\section{Part 10: Takeaway / Lessons Learned}
The primary takeaway is that user input should never be trusted, and development environments must never evaluate or execute code passed via external APIs. When handling serialized data like JSON, native and strictly typed parsers (like \texttt{JSON.parse()}) must be used instead of third-party libraries that rely on dynamic evaluation (like \texttt{eval()}). Furthermore, auditing dependencies for known vulnerabilities is critical for maintaining serverless security.
% \lessonchapter{02}{Broken Authentication}{authentication failure}
% \lessonchapter{03}{Sensitive Information Disclosure}{sensitive data exposure}
% \lessonchapter{04}{Insecure Cloud Configuration}{cloud misconfiguration}
% \lessonchapter{05}{Broken Access Control}{authorization failure}
% \lessonchapter{06}{Denial of Service}{availability failure}
% \lessonchapter{07}{Over-Privileged Function}{excessive IAM privilege}
% \lessonchapter{08}{Logic Vulnerabilities}{workflow or race-condition failure}
% \lessonchapter{09}{Vulnerable Dependencies}{dependency risk}
% \lessonchapter{10}{Unhandled Exceptions}{unsafe error handling}

\chapter*{Optional Bonus Findings}
\addcontentsline{toc}{chapter}{Optional Bonus Findings}
\placeholder{Use this section only if you discover valid additional findings beyond the 10 official lessons. Repeat the same reporting discipline if you include bonus work.}

\chapter*{Conclusion}
\addcontentsline{toc}{chapter}{Conclusion}
\placeholder{Summarize the project outcomes, the most important remediation themes, and the main security lessons learned from DVSA.}

\end{document}
